\chapter{博物馆}

\section{博物馆寄语汇编}


\section*{圆明园}

山峦不语，历史留音。

世界因多姿多彩的物种而精彩。每一个物种的消失，都是人类走向孤独的脚步。

\medskip

\noindent\textbf{圆明园：劫灰飞尽，笃行致远}

千年古都，西山脚下，一座圆明园，百年风霜路！一处大遗址，千载难逢时！党和政府的关怀、社会各界的关注使圆明园的灿烂文化得到阐释，同时，也使世人看到了一场文化劫难给国人带来的巨大震撼与伤痛。

看今朝，站在两个一百年的历史交汇点上，积淀深厚的圆明园在前所未有的发展形势下，各项文保工作成效卓著，遗址建设水平逐步提高。

展未来，圆明园将更好地发挥大遗址在传承守护优秀文化、开展爱国主义教育中的作用，激发爱国热情、凝聚人民力量、培育民族精神。为实现中华民族伟大复兴的中国梦做出积极贡献。

\newpage

\section*{苏州博物馆}

万年前，苏州从远古蒙昧中醒来。她风雨兼程，步履豪迈，走过茹毛饮血的石器时代，走过金戈铁马的春秋战国，见证了秦汉隋唐的天下一统，见证了宋元明清的锦绣繁华，饱经沧桑，历久而弥新，走进中华人民共和国这个崭新的伟大时代。

苏州如诗，是枫桥夜泊船，凌波过横塘；苏州如画，是姑苏台上月，人尽似神仙。且把吴钩看了，转瞬已是千年，但见水袖蹁跹，牙板轻拍，吴歙雅绝。天地悠悠，湖山点点，有一座城是苏州，是江南，既纯粹，又斑斓。烟波浩渺的太湖，是她的胸怀；奔流不息的长江，是她的血脉；山清水秀的江南，是她的舞台。立足中国，放眼世界，既古典，又现代，生机勃勃，方兴未艾。

万年的回眸，万年的等待，万年的苏州，继往开来。

\subsection*{成语及其解释}

\begin{enumerate}[leftmargin=*]
    \item \textbf{风雨兼程}：形容不顾艰难险阻，日夜兼程地赶路。比喻在困难和艰苦的环境中勇往直前。
    
    举例：为了按时完成项目，团队成员风雨兼程，不分昼夜地工作。
    
    \item \textbf{步履豪迈}：形容走路时步伐大而有力，也比喻做事有魄力，有气概。
    
    举例：他步履豪迈地走进会场，显得非常自信。
    
    \item \textbf{茹毛饮血}：原指原始社会人类不懂得用火，连毛带血地生吃禽兽。比喻处在未开化的野蛮状态。
    
    举例：在茹毛饮血的时代，人类的生存条件极其艰苦。
    
    \item \textbf{金戈铁马}：指战争，也形容战士持枪驰马的雄姿。宋代词人辛弃疾在《永遇乐·京口北固亭怀古》中有“想当年，金戈铁马，气吞万里如虎”的名句，表达了对往昔英雄气概的怀念和赞美。
    
    举例：历史上的金戈铁马，如今只能在史书和影视作品中窥见一斑。
    
    \item \textbf{天下一统}：指全国统一，也指思想、行动等方面的一致。
    
    举例：秦始皇统一六国，实现了天下一统。
    
    \item \textbf{锦绣繁华}：形容景色绚丽多彩，也形容城市繁荣昌盛。
    
    举例：这座城市的夜景锦绣繁华，灯火辉煌。
    
    \item \textbf{饱经沧桑}：指经历了许多世事变迁，也形容人经历了许多困苦和磨难。
    
    举例：这位老人饱经沧桑，脸上写满了岁月的痕迹。
    
    \item \textbf{历久弥新}：指经历长久的时间而更加显得新鲜，也形容某些事物或品质经得起时间的考验。
    
    举例：这件艺术品虽然年代久远，但历久弥新，依然受到人们的喜爱。
    
    \item \textbf{枫桥夜泊}：指唐代诗人张继的《枫桥夜泊》诗，后泛指夜泊枫桥。“月落乌啼霜满天，江枫渔火对愁眠。姑苏城外寒山寺，夜半钟声到客船。”这首诗描绘了诗人在枫桥边泊船时所见的夜景，与文中“枫桥夜泊船”的意象相呼应。
    
    举例：秋天的夜晚，独自漫步在枫桥边，仿佛能感受到“枫桥夜泊”的意境。
    
    \item \textbf{凌波微步}：形容女子步态轻盈，也比喻诗文、书法等轻快流畅。
    
    举例：她穿着高跟鞋，凌波微步地走进了宴会厅，吸引了所有人的目光。
    
    \item \textbf{姑苏台上月}：出自唐代诗人张继的《枫桥夜泊》，形容月色下的姑苏台（今苏州）景色美丽。“姑苏”是苏州的古称。
    
    举例：在姑苏台上月的照耀下，古城的夜晚显得格外宁静和迷人。
    
    \item \textbf{人尽似神仙}：形容人们生活得非常幸福，如同神仙一般。
    
    举例：这个小村庄的人们生活安宁，人尽似神仙。
    
    \item \textbf{吴钩}：古代吴地（今江苏一带）出产的一种弯刀，后泛指宝刀、利剑。在后世的文学作品中，“吴钩”也常被用来象征英勇、智慧和忠诚，成为英勇战士的代名词。李贺《南园十三首·其五》：“男儿何不带吴钩，收取关山五十州。请君暂上凌烟阁，若个书生万户侯？”
    
    举例：他收藏了一把古代的吴钩，锋利无比。
    
    \item \textbf{水袖蹁跹（pián xiān）}：形容舞者挥动长袖，轻盈起舞的样子。其中“蹁跹”一词被广泛使用，用以描绘美好的舞姿或者轻快的步伐。
    
    举例：舞台上的水袖蹁跹，舞姿曼妙，令人目不转睛。
    
    \item \textbf{牙板轻拍}：形容轻轻拍打节奏，常用于戏曲表演中。
    
    举例：戏曲演员在表演时，牙板轻拍，节奏感强烈。
    
    \item \textbf{吴歙雅绝}：形容吴地（今江苏一带）的风俗、文化等高雅独特。
    
    举例：苏州的园林建筑，体现了吴歙雅绝的风格。
    
    \item \textbf{天地悠悠}：形容天地辽阔，时间久远。
    
    举例：站在高山之巅，望着天地悠悠，不禁感慨万千。
    
    \item \textbf{湖山点点}：形容湖光山色点缀其间，景色宜人。
    
    举例：春天的西湖，湖山点点，美不胜收。
    
    \item \textbf{生机勃勃}：形容充满生机和活力。
    
    举例：春天的公园里，万物复苏，生机勃勃。
    
    \item \textbf{方兴未艾}：形容事物正在发展，尚未停止。
    
    举例：这个新兴行业方兴未艾，吸引了许多创业者的关注。
\end{enumerate}

\subsection*{历史时期注释}

\begin{enumerate}[leftmargin=*]
    \item \textbf{石器时代}：指的是人类使用石制工具的时期，大约从距今约260万年前开始，直到约5000年前结束。这个时期包括旧石器时代和新石器时代。
    
    \item \textbf{春秋战国}：指的是中国历史上春秋（公元前770年—公元前476年）和战国（公元前475年—公元前221年）两个时期，这一时期以诸侯争霸和诸子百家的思想争鸣为特点。
    
    \item \textbf{秦汉隋唐}：指的是中国历史上的秦朝（公元前221年—公元前206年）、汉朝（公元前202年—公元220年，包括西汉和东汉）、隋朝（公元581年—公元618年）和唐朝（公元618年—公元907年）。这些朝代见证了中国历史上的多次统一和中央集权的加强。
    
    \item \textbf{宋元明清}：指的是中国历史上的宋朝（公元960年—公元1279年，包括北宋和南宋）、元朝（公元1271年—公元1368年）、明朝（公元1368年—公元1644年）和清朝（公元1636年/1644年—公元1912年）。这些朝代时期，中国社会经济和文化发展到了一个新的高度，尤其是宋朝，被认为是中国历史上的文艺复兴时期。
\end{enumerate}

\newpage

\section*{西安博物馆：让过去拥有未来}

很多年以后，蓝衣老人还清晰记得自己初入长安那种直击心底的震撼。那时，他还是一个蓝衣少年，从祖辈、父辈那里听得太多太多关于长安的故事。

他们说，那里诗礼昌盛、艺术璀璨，千门万户、四通八达，店肆林立、百业兴旺，胡风盛行、百戏喧腾，俯拾即是机遇。

长安梦深深植根于心底，怀抱热望，少年踏上了逐梦之旅。黄沙漫漫，驼铃声声。

长安，越来越近……

\subsection*{词语注释}

\begin{itemize}[leftmargin=*]
    \item \textbf{店肆林立}：形容商店很多，一个接一个地排列着。
    \item \textbf{诗礼昌盛}：指诗歌和礼仪都非常繁荣兴盛。
    \item \textbf{俯拾即是}：形容某种东西非常常见，很容易得到。
\end{itemize}

\newpage

\section*{成都武侯祠}

这是一个英雄辈出，逆流而上的时代。

日月旋转的宏图野望，沧海横流的英雄豪情，那么多闪烁着智慧与生命之美的个体，在星汉灿烂之下，播扬激荡昭彰的人生理想，逐鹿中原，酝酿着权力分合的秩序重构。在这段洪波涌动的历史河流中，刘备与诸葛亮二人相逢于乱世，风云际会，托付以家国远志。君明臣良，厚相结纳，忠信节义，一时无双，于临危之际信义相托，于险阻之中肝胆相照。浮光掠影间麈尾轻扬，揽江山，听风雨，他们共叙经纬，他们挥剑指北，他们山水跋涉，他们长揖作别。明月昭昭，君臣义，长风起兮，鱼水契。

\subsection*{成语及其解释}

\begin{enumerate}[leftmargin=*]
    \item \textbf{日月旋转}：形容时间的流逝，岁月的更替。这个成语用来形容时间的快速流逝，如同日月在天空中不断旋转。
    
    举例：历史的车轮日月旋转，见证了无数王朝的兴衰更迭。
    
    \item \textbf{沧海横流}：形容社会动荡不安，局势混乱。这个成语原本指大海波涛汹涌，后来用来比喻社会动荡，局势不稳定。
    
    举例：在那个沧海横流的年代，英雄豪杰纷纷崛起，试图力挽狂澜。
    
    \item \textbf{星汉灿烂}：形容星空璀璨，也比喻人才荟萃。
    
    举例：在那个星汉灿烂的时代，无数文人墨客竞相争辉，留下了许多传世佳作。
    
    \item \textbf{播扬激荡}：形容传播和激荡，使影响广泛。
    
    举例：他的演讲播扬激荡，激发了听众的热情和斗志。
    
    \item \textbf{逐鹿中原}：比喻争夺天下或权力。《史记·淮阴侯列传》：“秦失其鹿，天下共逐之。”这句话是成语“逐鹿中原”的出处，形象地描述了秦朝末年天下大乱，群雄并起争夺天下的情形。刘伯承《千里跃进大别山》：“三军在江、淮、河、汉之间布成‘品’字形阵势，互为犄角，逐鹿中原，机动歼敌。”
    
    举例：在那个群雄逐鹿中原的时代，只有智勇双全者才能最终胜出。
    
    \item \textbf{风云际会}：形容英雄豪杰在动荡的时代相遇，共同成就大业。
    
    举例：在那个风云际会的时代，许多英雄人物汇聚一堂，共同推动了历史的进程。
    
    \item \textbf{君明臣良}：形容君主英明，臣子贤良。
    
    举例：在那个君明臣良的时代，国家政治清明，百姓安居乐业。
    
    \item \textbf{忠信节义}：形容忠诚、诚信、有节操和正义。
    
    举例：他一生忠信节义，深受人们的尊敬和爱戴。
    
    \item \textbf{肝胆相照}：形容彼此之间非常坦诚，真心相待。《史记·淮阴侯列传》：“臣愿披腹心，输肝胆，效愚计，恐足下不能用也。”这句话是成语“肝胆相照”的出处，描述了蒯通愿意向韩信展示自己的忠诚和诚意。
    
    举例：他们两人肝胆相照，共同经历了许多艰难险阻。
    
    \item \textbf{浮光掠影}：形容事物留下的印象不深，一晃而过。清代李汝珍所著《镜花缘》中有一句：“学问从实地上用功，议论自然确有根据；若浮光掠影，中无成见，自然随波逐流，无所适从。”
    
    举例：对于那段历史，他只有浮光掠影的记忆，难以形成完整的认识。
\end{enumerate}

\newpage

\section*{上海博物馆}

从长三角腹地的福泉山、吴淞江畔的青龙镇到黄浦江边的上海城，上海的历史波澜跌宕，从不乏激动人心的闪耀时刻。考古学家们通过对实物遗存的艰苦发掘与不懈研究，将无数的信息碎片连缀成历史，还原那些风云激荡的精彩故事，完成对上海历史的考古式呈现。

从考古的视角回望，是六千年的深厚积淀塑造了今天的上海和生活在这座城市的人们。水与土，构筑沃野良港，成就了海纳百川的宽广胸怀；城与人，彼此相映成辉，奠定了开明睿智的城市精神。考古，揭示出我们曾经的来路；历史，将启迪我们今天与未来的新途。

上海，文脉璀璨；上海，繁华依然。

\newpage

\section*{信阳市博物馆}

一花一世界，一叶一菩提。一片片浮雕的拼合，一件件造像的归复，无不在提醒你我这样的见微知著者，“国运弱则文化难，国运强则文化盛”。

昔日，盗凿劫掠，流散异乡，文化之殇，民族之痛！今朝，回归故里，聚合复位，文化之兴，民族之幸！

以物记事，以事叙史，以史启思。在龙门有窟无数的残首断臂中回首历史的沧桑，在一件件身首合璧的造像中展望更多的国宝回家。

兴文化、强国运。道阻且长，行则将至；行而不辍，未来可期。

\subsection*{文化注释}

\textbf{“一花一世界，一叶一菩提”} 源于《华严经》，其原文是：“佛土生五色茎，一花一世界，一叶一如来。”这句表达了华严境界，本义是说：微尘里面有世界，世界里面有微尘，重重无尽，主伴相融，相摄相入，无碍自在。这句话蕴含着深刻的哲学意味，表达的是即使是微小的事物也包含着整个世界的真理和智慧，体现了佛教中微观与宏观相互关联、一即一切、一切即一的宇宙观。

\subsection*{成语解释}

\textbf{见微知著}：意指从微小的苗头或迹象中，就能预见到事物的发展趋势和结果。这个成语强调通过观察细节来洞察整体，从小处着眼来理解大的格局。

应用示例：
\begin{itemize}[leftmargin=*]
    \item 天气预报：气象学家通过观察天空中云彩的微小变化，预测即将到来的天气变化。
    \item 医学诊断：医生通过观察患者的细微症状，如体温、血压和皮肤颜色，来诊断疾病。
    \item 商业趋势：市场分析师通过观察消费者的微小购买习惯，预测市场趋势和消费者需求。
    \item 社交动态：通过观察一个人的微小行为，了解他们的情绪、态度和意图。
    \item 科学发现：科学家通过仔细观察和实验，发现微小的物理现象或化学变化，推动科学的进步。
    \item 环境保护：环保工作者通过观察环境中的微小变化，评估环境健康状况。
    \item 犯罪侦查：侦探通过分析现场的微小线索，如指纹、毛发或纤维，追踪犯罪嫌疑人和破案。
\end{itemize}

\subsection*{词语释义}

\textbf{文化之殇}：“文化之殇”蕴含着深刻的悲伤、痛苦或失落感，特别是在文化领域中的重大损失或断裂。

\begin{enumerate}[leftmargin=*]
    \item \textbf{深刻的悲伤与失落}：用来表达对某种文化价值、文化遗产或文化传统的丧失所感受到的深切悲痛。
    \item \textbf{文化断层与遗忘}：城市化和现代化进程中，历史遗迹、文化古迹和人文底蕴因缺乏保护而遭到破坏，导致文化功能的缺失和文明品格的不健全。
    \item \textbf{历史耻辱与文化自卑}：中国近代史上，由于西方列强的侵略和压迫，中国的海洋文化受到屈辱，这种历史之辱被视为中国近代文化之殇的起点。
    \item \textbf{文化冲突与解构}：传统文化在西方文化和世俗文化的冲击下所遭受的解构和溃败。
    \item \textbf{文化自我认同的丧失}：全球化影响下，本土文化被外来文化取代，导致民族文化的边缘化和自我认同的丧失。
    \item \textbf{文化符号的丧失}：一些文化符号和传统逐渐失去它们原有的意义和色彩。
\end{enumerate}

\textbf{殇}的含义：
\begin{enumerate}[leftmargin=*]
    \item \textbf{未成年而死}：未达到成年年龄就去世的人。古代文献中有“长殇”、“中殇”、“下殇”等称谓。
    \item \textbf{为国战死者}：即“国殇”，用以表达对战争中牺牲者的哀悼和尊敬。
    \item \textbf{横死、非正常死亡}：意外或暴力导致的死亡。
\end{enumerate}

\newpage

\section*{新疆博物馆}

天山昆仑、大漠戈壁，“大美新疆”不仅有壮阔美景，更有令人惊叹的古代遗存，向世人娓娓叙说中华文明的绵远与博大。

残垣断壁，尺砖片瓦，无不见证了新疆是中国不可分割的一部分，见证了新疆各民族是中华民族血脉相连的家庭成员，亲历了包括新疆各族人民在内的全体中华儿女共同奋斗创造悠久灿烂中华文明的伟大历程，更目睹了古今丝绸之路的繁盛。从落后、被动，到自主、开拓，再到自觉与创新，新疆考古已经走过不平凡的百年。

日月交替，不变的是考古人赓续中华文脉、大书中华文明的赤诚之心。学术探索是永无止境的求真之旅，上穷碧落下黄泉，再现文明的光辉，守正创新，行稳致远。新疆考古在弘扬文化自信的道路上将书写出更加浓墨重彩的一笔。中国故事是永恒的话题，新疆篇章精彩动人。

\newpage

\section*{广州博物馆}

唐宋时期，海上丝绸之路进入繁荣时期。广州是举世闻名的东方大港、“通海夷道”的东方起点，东西方商品在这里聚集并转口输往海内外市场。隋朝，在广州设立南海神庙，为历代尊崇。唐朝，在广州首设市舶使（院），开创古代市舶管理制度，掌管海外贸易。宋朝，制订“广州市舶条”，将市舶制度推广到国内其他港口。经海路，中国与国外的官方与民间交流得以加强，远赴印度的佛教求法僧、在广州“蕃坊”长期生活的各国商人都是鲜活的例证。唐宋时期广州与海上丝绸之路是一部跨越“七海”的商贸史和交流史，是当今建设21世纪海上丝绸之路的历史见证和文化源泉。

\newpage

\section*{国家自然博物馆·恐龙部分结语}

如果把地球的历史比作一天的话，凌晨0点地球形成，现在是24点，那么恐龙22:48出现，23:40灭亡，它们在地球上生活了52分钟。

而我们人类即使从周口店北京猿人算起也只有9.4秒钟。

和恐龙比起来，我们在地球上的时间太短了。可是现在已经面临着许多生存问题，所以我们一定要珍惜环境，珍爱生命。

\newpage

\section*{附录：《赤壁赋》节选}

作者：苏轼

\medskip

\noindent 况吾与子渔樵于江渚之上，侣鱼虾而友麋鹿，驾一叶之扁舟，举匏樽以相属。寄蜉蝣于天地，渺沧海之一粟。

\subsection*{注释}

\begin{itemize}[leftmargin=*]
    \item \textbf{渔樵（yú qiáo）}：指捕鱼和砍柴，这里代指隐居生活。
    \item \textbf{江渚（jiāng zhǔ）}：江中的小洲或岸边。
    \item \textbf{侣（lǚ）}：以……为伴侣，名词的意动用法。
    \item \textbf{麋鹿（mí lù）}：麋，鹿的一种，这里泛指鹿。
    \item \textbf{扁舟（piān zhōu）}：小船。
    \item \textbf{匏樽（páo zūn）}：匏制的酒樽，亦泛指饮具。
    \item \textbf{蜉蝣（fú yóu）}：一种昆虫，夏秋之交生于水边，生命短暂，仅数小时。此句比喻人生之短暂。
    \item \textbf{渺（miǎo）}：小。
    \item \textbf{沧海（cāng hǎi）}：大海。
    \item \textbf{粟（sù）}：谷子，这里比喻极小的东西。
\end{itemize}
